% In this file you should put the actual content of the blueprint.
% It will be used both by the web and the print version.
% It should *not* include the \begin{document}
%
% If you want to split the blueprint content into several files then
% the current file can be a simple sequence of \input. Otherwise It
% can start with a \section or \chapter for instance.

\chapter{Current library}

\section{\texttt{int\_seq}}

This library file proves that certain native datatypes, such as \texttt{uint63}, and their operations
reflect their mathematical counterparts. These allows for computation speedup in rest of the database.

\section{\texttt{sizelexi}}

This library file establishes the well-foundedness of the lenlex order.
For the remainder of the section we fix an alphabet $A$ and an order $\leq$ on $A$.

\begin{definition}[Lenlex order]%
	\label{def:sizelexi}
	\rocq{MonPres.sizelexi.sizelexi}
	\rocqok{}
	The \emph{lenlex order} $\lenlex$ induced by $(A, \leq)$ is the relation on $A^\ast$ given by
	\[ u \lenlex v \quad \text{if} \quad |u| < |v| \quad \text{or} \quad \left(|u| = |v| \text{ and } u\lex v\right), \]
	where $\lex$ is the usual lexicographic order induced by $A$ under $\leq$.
\end{definition}

\begin{lemma}%
	\label{lem:sizelexi_total_order}
	\rocqok{}
	\uses{def:sizelexi}
	If $\leq$ is total, then the lenlex order induced by $(A, \leq)$ is a total order and $\varepsilon$ is its minimum
	element.
\end{lemma}
\begin{proof}
	\rocqok{}
\end{proof}

\begin{lemma}%
	\label{lem:lt_sizelexi_stable}
	\rocq{MonPres.sizelexi.lt_sizelexi_stable}
	\rocqok{}
	\uses{def:sizelexi}
	The lenlex order induced by $(A, \leq)$ is stable under right and left multiplication by elements of $A^\ast$. That is,
	\[u\lenlex v \quad \Rightarrow \quad sut \lenlex svt \]
	for all $s, t\in A^\ast$.
\end{lemma}
\begin{proof}
	\rocqok{}
\end{proof}

\begin{lemma}%
	\label{lem:sizelexi_wf}
	\rocq{MonPres.sizelexi.sizelexi_wf}
	\rocqok{}
	\uses{def:sizelexi}
	If $\leq$ is well founded, then the lenlex order induced by $(A, \leq)$ under $\leq$ is also well founded. That is,
	every subset of $A^\ast$ admits a minimum element.
\end{lemma}
\begin{proof}
	\rocqok{}
\end{proof}

\section{\texttt{monoids}}

This file establishes some important background regarding the theory of monoids.

\begin{definition}%
	\label{def:FreeMonoidT}
	\rocq{MonPres.monoids.FreeMonoidT}
	\rocqok{}
	Let $A$ be a set. The \emph{free monoid} generated by $A$ is the monoid $A^\ast$ whose elements are words over the
	alphabet $A$ with multiplication given by concatenation. The identity is the empty word $\varepsilon$.
\end{definition}

\begin{definition}%
	\label{def:univmor}
	\rocq{MonPres.monoids.univmor}
	\rocqok{}
	\uses{def:FreeMonoidT}
	Let $A$ be a set, let $M$ be a monoid and let $f: A \rightarrow M$ be a function.
	Define the \emph{universal morphism} from $A^\ast$ into $M$ induced by $f$ to be the function
	$\hat{f}: A^\ast \rightarrow M$ given by
	\[(a_1\cdots a_n)\hat{f} = (a_1)f\cdots(a_n)f\]
	for all $n\in \N_0$ and $a_1, \ldots, a_n\in A$.
\end{definition}

\begin{lemma}%
	\label{lem:univmor_uniq}
	\rocq{MonPres.monoids.univmor_uniq}
	\rocqok{}
	\uses{def:univmor}
	Let $A$ be a set, let $M$ be a monoid and let $f: A \rightarrow M$ be a function.
	Then $\hat{f}$ is the unique monoid homomorphism $A \rightarrow M$ such that $(a)\hat{f} = (a)f$ for all $a\in A$.
\end{lemma}
\begin{proof}
	\rocqok{}
\end{proof}

\chapter{Preliminaries}\label{chap:preliminaries}

In this chapter we establish the basic theoretic underpinnings of monoid theory, with a focus on monoid and group
presentations and the word problem.

\section{Finitely presented monoids}

\begin{definition}\label{def:two_sided_cong}
	\rocq{MonPres.present.congruencep}
	\rocqok{}
	Let $M$ be a monoid and let $\rho \subseteq M\times M$. We say that $\rho$ is an \emph{two-sided congruence} of $M$
	if $\rho$ is an equivalence relation such that $(x, y)\in \rho \Rightarrow (sxt, syt)\in \rho$ for all
	$x, y, s, t\in M$.
\end{definition}

\begin{definition}\label{def:equal_mod_R}
	\catcode`\%=11\relax
	\rocq{MonPres.present.Defs%2EDefPresentationRels%2Eequiv_to
	}
	\catcode`\%=14\relax
	% Note: above is a very bad hack to make the linking work. It will be fixed in a future version of
	% rocqblueprint.
	\rocqok{}
	Let $A$ be a set, let $R\subseteq A^\ast \times A^\ast$ and let $u, v\in A^\ast$. We say that $u$ and $v$ are
	\emph{equal modulo} $R$, written $u =_R v$, if there exists an $n\in \N$ and sequences $w_1, \ldots, w_{n+1}$,
	$s_1, \ldots, s_n$, $t_1, \ldots, t_n$, $u_1, \ldots, u_n$, $v_1, \ldots, v_n\in A^\ast$ such that
	\begin{enumerate}
		\item $u = w_1$,
		\item $v = w_{n+1}$,
		\item for all $i\in \{1, \ldots, n\}$, $s_i u_i t_i = w_i$,
		\item for all $i\in \{1, \ldots, n\}$, $s_i v_i t_i = w_{i+1}$,
		\item for all $i\in \{1, \ldots, n\}$, either $(u_i, v_i) \in R$ or $(v_i, u_i) \in R$.
	\end{enumerate}
\end{definition}

\begin{lemma}\label{lem:equal_mod_R_cong}
	\rocq{MonPres.present.equiv_congr}
	\rocqok{}
	\uses{def:equal_mod_R}
	Let $A$ be a set, let $R\subseteq A^\ast \times A^\ast$. Then $=_R$ is a two-sided congruence of $A^\ast$.
\end{lemma}
\begin{proof}
	\rocqok{}
\end{proof}

\begin{lemma}\label{lem:equal_mod_R_is_cong_gen_by_R}
	\rocq{MonPres.present.equiv_min}
	\rocqok{}
	\uses{def:equal_mod_R, def:two_sided_cong}
	Let $A$ be a set, let $R\subseteq A^\ast \times A^\ast$ and let $\rho$ be a two-sided congruence of $A^\ast$.
	If $R \subseteq \rho$, then $=_R \subseteq \rho$.
\end{lemma}
\begin{proof}
	\uses{lem:equal_mod_R_cong}
	\rocqok{}
	If $(u, v)\in R$, then $u=_R v$ follows by taking $n = 1$, $w_1 = u_1 = u, w_2 = v_1 = v$ and
	$s_1 = t_1 = \varepsilon$. Hence $=_R$ contains $R$. By \Cref{lem:equal_mod_R_cong}, $=_R$ is also a two-sided
	congruence. If we can show that $=_R$ is contained in $R^\#$, then the minimality of $R^\#$ will force $=_R$ to
	coincide with $R^\#$. Let $u, v\in A^\ast$ be any pair such that $u =_R v$ holds and let $w_i, u_i, v_i, s_i, t_i$ be
	the associated sequences of words. It follows that
	$(u_i, v_i)\in R^\#$ for all $i\in\{1,\ldots, n\}$, since either $(u_i, v_i)\in R$ or $(v_i, u_i)\in R$ and
	$R^\#$ is a symmetric relation containing $R$. Furthermore, $(w_i, w_{i+1}) = (s_i u_i t_i, s_i v_i t_i)\in R^\#$ for
	all $i\in \{1, \ldots, n\}$
	since $R^\#$ is a two-sided congruence. Finally, $(u, v) = (w_1, w_{n+1}) \in R^\#$ by transitivity of $R^\#$. Hence
	$=_R$ is contained in $R^\#$ as required.
\end{proof}

\begin{definition}\label{def:monoid_presentation}
	Let $M$ be a monoid, let $A$ be a set and let $R\subseteq A^\ast \times A^\ast$. We say
	that $\Mon\presset{A}{R}$ is a \emph{monoid presentation} for $M$ if $M \cong A^\ast / R^\#$. In this case we say that
	$A$ is an \emph{generating set} for $M$ and that $R$ is a \emph{set of relations}.
\end{definition}

\begin{definition}\label{def:generating_set_inclusion}
	\uses{def:monoid_presentation}
	Let $M$ be a monoid presented by $\Mon\presset{A}{R}$ and let $\psi: A^\ast/R^\# \rightarrow M$ be an isomorphism.
	The \emph{inclusion} of $A$ in $M$ with respect to $\psi$ is the map $\iota: A \rightarrow M$ given by taking
	$(a)\iota = (a/R^\#)\psi$. Due to the universal property of $A^\ast$, every inclusion of $A$ in $M$ extends uniquely
	to a homomorphism $\varphi: A^\ast\rightarrow M$ given by $(u)\varphi = (u/R^\#)/\psi$.
\end{definition}

\begin{lemma}\label{lem:equal_mod_R_is_equal_phi}
	\uses{def:generating_set_inclusion, def:equal_mod_R, def:monoid_presentation}
	Let $M$ be a monoid presented by $\Mon\presset{A}{R}$ and let $\varphi: A^\ast\rightarrow M$ be a monoid homomorphism
	extending an inclusion of $A$ in $M$. Then $u=_R v$ if and only if $(u)\varphi = (v)\varphi$.
\end{lemma}
\begin{proof}
	\uses{lem:equal_mod_R_is_cong_gen_by_R}
	Let $\psi: A^\ast / R^\# \rightarrow M$ be the isomorphism such that $(u/R^\#)\psi = (u)\varphi$ for all
	$u\in A^\ast$. Then $(u)\varphi = (v)\varphi$ if and only if $(u/R^\#)\psi = (v/R^\#)\psi$. Since $\psi$ is an
	isomorphism, the latter occurs precisely when $u/R^\# = v/R^\#$, or, equivalently $(u, v)\in R^\#$. It follows by
	\Cref{lem:equal_mod_R_is_cong_gen_by_R} that $(u, v)\in R^\#$ if and only if $u=_R v$. This concludes the proof.
\end{proof}

\section{Finitely presented groups}

\begin{definition}\label{def:reduced_word}
	Let $A$ be a set, let $A^{-1}$ be a set of formal inverses of $A$ and let $w\in {\left(A\cup A^{-1}\right)}^\ast$.
	We say that $w$ is \emph{freely reduced} if $w$ does not contain $aa^{-1}$ or $a^{-1}a$ as a subword for all $a\in A$.
\end{definition}

\begin{definition}\label{def:group_presentation}
	\uses{def:monoid_presentation}
	Let $G$ be a group, let $A$ be a set, let $A^{-1}$ be a set of formal inverses of $A$ and let
	$R\subseteq {\left(A\cup A^{-1}\right)}^\ast$. We say that $G$ is presented as a group by $\Gp\presset{A}{R}$ if it is
	presented as a monoid by
	\[
		\Mon\presset{A\cup A^{-1}}{R\cup\set{aa^{-1} = 1}{a\in A}\cup\set{a^{-1}a=1}{a\in A}}.
	\]
\end{definition}

\begin{definition}\label{def:free_group}
	\uses{def:group_presentation}
	Let $A$ be a set. The \emph{free group} with alphabet $A$ is the group $F_A$ presented by
	$\Gp\presset{A}{\varnothing}$.
\end{definition}

\begin{lemma}[Universal property of the free group]\label{lem:universal_prop_free_group}
	\uses{def:free_group}
	Let $A$ be a set, let $F_A$ be the free group with alphabet $A$ and let $\iota: A\rightarrow F_A$ be the trivial
	inclusion. Then for every group $G$ and every map $f: A\rightarrow G$ there exists a unique group homomorphism
	$\varphi: F_A \rightarrow G$ such that $f = \iota \circ \varphi$.
\end{lemma}

\begin{definition}\label{def:free_basis}
	\uses{def:free_group, lem:universal_prop_free_group}
	Let $G$ be a group, let $A\subseteq G$ be any subset and let $\varphi: F_A \rightarrow G$ be the group homomorphism
	extending the inclusion of $A$ in $G$. We say that $A$ is a \emph{basis of a free subgroup} of $G$ if $\varphi$ is
	injective.
\end{definition}

\section{The word problem}

\begin{definition}\label{def:WPdecidable}
	\uses{def:monoid_presentation}
	\rocq{MonPres.present.WPdecidable}
	\rocqok{}
	% TODO(reiniscirpons): I am not so sure this is the correct definition we want. It is not wrong per-se, but
	% it might be preferable to define decidability of word problem in terms of the equality relation of the monoid.
	Let $M$ be a monoid presented by the monoid presentation $\Mon\presset{A}{R}$. We say that $M$ has
	\emph{decidable word problem} if for all $u, v\in A^\ast$ it is decidable if $u =_R v$.
\end{definition}

\begin{lemma}\label{lem:embed_preimage_preserves_dec}
	\uses{def:WPdecidable}
	If a monoid $M$ embeds into a monoid $M'$ and $M'$ has decidable word problem, then $M$ also has
	decidable word problem.
\end{lemma}
\begin{proof}
	% NOTE(reiniscirpons): This proof is what bites us because of the awkward definition above.
	\uses{lem:equal_mod_R_is_equal_phi}
	Let $\Mon\presset{A}{R}$ be a monoid presentation for $M$ and let $\Mon\presset{A^\prime}{R^\prime}$ be a presentation
	for $M^\prime$. Let $\varphi: A^\ast\rightarrow M$ be a monoid homomorphism extending an inclusion of $A$ in
	$M$, let $\varphi^\prime: A^{\prime\ast} \rightarrow M^\prime$ be defined analogously for $M^\prime$ and let
	$\iota: M \rightarrow M'$ be an embedding.

	Take $u, v\in A^\ast$ to be arbitrary. Since $A^\prime$ is a generating set for $M^\prime$, there exist
	$u^\prime, v^\prime \in A^{\prime\ast}$ such that $(u^\prime)\varphi^\prime = (u)\varphi\circ \iota$ and
	$(v^\prime)\varphi^\prime = (v)\varphi\circ\iota$. Assume that $u^\prime =_{R^\prime} v^\prime$ in $M^\prime$, then
	$(u^\prime)\varphi^\prime = (v^\prime)\varphi^\prime$ by \Cref{lem:equal_mod_R_is_equal_phi}.
	Since $\iota$ is injective it follows that $(u)\varphi = (v)\varphi$. By \Cref{lem:equal_mod_R_is_equal_phi}
	again, $u=_R v$. So $u^\prime =_{R^\prime} v^\prime$ implies $u=_R v$.

	Now assume $u=_R v$ instead. Then, reasoning as before, $(u)\varphi = (v)\varphi$, hence
	\[
		(v^\prime)\varphi^\prime = (v)\varphi\circ\iota =
		(u)\varphi\circ\iota = (u^\prime)\varphi^\prime.
	\]
	And by a final application of \Cref{lem:equal_mod_R_is_equal_phi}, we have that $u^\prime =_{R^\prime} v^\prime$.

	Hence for every pair of words $u, v\in A^\ast$ there exists a pair of words $u^\prime, v^\prime \in A^\ast$ whose
	equality modulo $R^\prime$ is equivalent to the euqality modulo $R$ of $u$ and $v$. Hence we can use the decision
	procedure deciding the word problem in $M^\prime$ to decide the word problem in $M$.
\end{proof}

\chapter{The word problem in 1-relation groups}\label{chap:1_relation_groups}

In this chapter we establish the decidability of the word problem for $1$-relation groups in
\Cref{thm:group_1rel_dec}. The proofs that special $1$-relation monoids and group embeddable $1$-relation monoids
have decidable word problem both reduce to the group case, so we take care of it here.

\section{HNN-extensions}

\begin{definition}\label{def:hnn_extension}
	Let $G$ be a group presented by $\Gp\presset{A}{R}$ and let $H, K$ be isomorphic subgroups of $G$.
	The \emph{HNN-extension} of $G$ with respect to the isomorphism $\varphi: H \rightarrow K$ is the group
	$G\ast_\varphi$ presented by
	\[\Gp\presset{A\cup \{t\}}{R\cup\set{(t h t^{-1}, (h)\varphi)}{h\in H}},\]
	where $t$ is a new symbol not in $A$.
\end{definition}

\section{The Freiheitssatz and the Magnus-Moldavanskii hierarchy}

Our primary references for this section are~\cite{Moldavanskii1967aa,McCool1973aa,Linton2025aa} and we sometimes
follow~\cite{PutnamYYYYaa} in terms of exposition and phrasing.

\begin{lemma}[c.f.~{\cite[Lemma 2.2]{PutnamYYYYaa}}]\label{lem:1rel_group_embedding}
	\uses{def:group_presentation}
	Let $G$ be a group presented by $\Gp\presset{A}{r = 1}$, let $a, b\in A$ be distinct elements and let
	$\alpha, \beta\in \Z$ be non-zero.
	Let $\psi: F_A \rightarrow F_A$ be the unique homomorphism such that
	\[
		(c)\psi = \begin{cases}
			a^\beta   & \text{if }c=a     \\
			ba^\alpha & \text{if }c=b     \\
			c         & \text{otherwise}.
		\end{cases}
	\]
	Let $G^\prime$ be the group presented by $\Gp\presset{A}{(r)\psi = 1}$, let $\varphi:F_A\rightarrow G$ and
	$\varphi^\prime: F_A\rightarrow G^\prime$ be the usual maps extending the inclusion of $A$ in $G$ and $A$ in
	$G^\prime$, respectively. Finally, let $\overline{\psi}: G\rightarrow G^\prime$ be the map given by
	\[((w)\varphi)\overline{\psi} = (w)\psi \circ \varphi^\prime\]
	for all $w\in F_A$. Then $\overline{\psi}$ is a well defined injective group homomorphism.
\end{lemma}

\begin{lemma}\label{lem:Magnus_induction_principle}
	\uses{def:free_group, def:hnn_extension}
	% TODO(reiniscirpons): This might not be actually correct
	Let $\mathscr{P}$ be a property that holds for all free products of a free group and a cyclic group and is preserved
	under taking preimages of embeddings and HNN-extensions over free sugroups. Then $\mathscr{P}$ holds for all
	1-relation groups.
\end{lemma}
\begin{proof}
	\uses{lem:1rel_group_embedding}
\end{proof}


\begin{lemma}[Freiheitssatz, c.f.~{\cite{Magnus1932aa}}]\label{lem:Freiheitssatz}
	\uses{def:WPdecidable, def:group_presentation, def:reduced_word, def:free_basis}
	Let $G$ be a group presented by $\Gp\presset{A}{r = 1}$, where $r$ is a freely reduced word, and let $B\subseteq A$.
	If $r$ cannot be written as a word in $B\cup B^{-1}$, then $B$ is a basis of a free subgroup of $G$.
\end{lemma}
\begin{proof}
	\uses{lem:Magnus_induction_principle}
\end{proof}


\begin{theorem}[c.f.~{\cite{Magnus1932aa}}]\label{thm:group_1rel_dec}
	\uses{def:WPdecidable, def:group_presentation}
	Every group presented by a $1$-relation group presentation has decidable word problem.
\end{theorem}
\begin{proof}
	\uses{lem:Freiheitssatz, lem:Magnus_induction_principle}
\end{proof}

\chapter{Special 1-relation monoids}\label{chap:special_1_rel_monoids}

\chapter{Group embeddings and Adian's algorithm}\label{chap:group_embeddings_and_adians_algorithm}

\begin{theorem}[c.f.~{\cite[Theorem 4]{Adian1960aa}}]\label{thm:cycle_free_embeds_group}
	\uses{def:WPdecidable}
	Every cycle-free monoid embeds in a group. Furthermore, if $\Mon\presset{A}{R}$ is a presentation for the cycle free
	monoid, then $\Gp\presset{A}{R}$ is a presentation for a group into which it embeds, under the inclusion embedding
	$(a)\iota = a$ for all $a\in A$.
\end{theorem}

\begin{theorem}\label{thm:cycle_free_1rel_dec}
	\rocq{MonPres.criteria.cycle_free_1rel_dec}
	\rocqok{}
	\uses{def:WPdecidable}
	If a $1$-relation monoid presentation is cycle-free, then the underlying monoid has decidable word problem.
\end{theorem}
\begin{proof}
	\uses{thm:cycle_free_embeds_group, thm:group_1rel_dec, lem:embed_preimage_preserves_dec}
	By \Cref{thm:cycle_free_embeds_group}, every cycle-free monoid embeds in the group with the same presentation.
	By \Cref{thm:group_1rel_dec}, every 1-relation group has decidable word problem. Finally,
	\Cref{lem:embed_preimage_preserves_dec}
	allows us to use the decidability in the embedding to conclude decidability of the embedded group.
\end{proof}

\chapter{Small overlap monoids}\label{chap:small_overlap}

\bibliographystyle{tugboat}
\bibliography{biblio}
